
\usepackage[T1]{fontenc}
\usepackage{lmodern}
\usepackage[margin=1in]{geometry}

\usepackage{xparse}
\usepackage{ifthen}

\usepackage{adjustbox}
\usepackage{changepage}
\usepackage{comment}
\usepackage{enumitem}
\usepackage{graphicx}
\usepackage{lastpage}
\usepackage{needspace}
\usepackage{xcolor}

\usepackage{amsmath}
\usepackage{amssymb}
\usepackage{amsthm}
\usepackage{mathtools}
\usepackage{siunitx}

\usepackage{pdfpages}
\usepackage{tocloft}

\usepackage{tikz}
\usepackage{pgfplots}
\pgfplotsset{compat=1.18}
\usetikzlibrary{arrows.meta, positioning, shapes.geometric}
\usetikzlibrary{datavisualization, datavisualization.formats.functions}


%%%%% Hyperlinks
\usepackage{hyperref}
\hypersetup{
	colorlinks=true,
	urlcolor=blue,
	linkcolor=blue,
	citecolor=red
}

%%%%% Math & Table Configuration
\sisetup{
	table-number-alignment=center,
	mode=math
}

%%%%% Custom Math Commands
\newcommand{\mbb}[1]{\mathbb{#1}}
\newcommand{\mrm}[1]{\mathrm{#1}}
\newcommand{\mcal}[1]{\mathcal{#1}}
\newcommand{\msf}[1]{\mathsf{#1}}
\newcommand{\mf}[1]{\mathfrak{#1}}
\newcommand{\mscr}[1]{\mathscr{#1}}

%%%%% Symbols
\newcommand{\abs}[1]{\mathopen{}\lvert #1 \rvert\mathopen{}}
\newcommand{\lrp}[1]{\mathopen{}\left(#1\right)\mathclose{}}
\newcommand{\lrb}[1]{\mathopen{}\left[#1\right]\mathclose{}}
\newcommand{\lrcb}[1]{\mathopen{}\left\{#1\right\}\mathopen{}}
\newcommand{\lrsb}[1]{\mathopen{}\left[#1\right]\mathopen{}}
\newcommand{\dd}{\mrm{d}}
\NewDocumentCommand{\deriv}{O{} m m}{\frac{\dd^{#1} \IfValueTF{#2}{#2}{}}{\dd #3^{#1}}}
\NewDocumentCommand{\inte}{O{} O{} m m}{\int_{#1}^{#2} #3\, \dd #4}
\newcommand{\todo}{\textcolor{red}{TODO}}

%%%%% Theorem Environments
\newtheorem*{thm*}{Theorem}
\newtheorem*{defn*}{Definition}
\newtheorem*{cor*}{Corollary}
\newtheorem*{prop*}{Proposition}
\newtheorem*{lem*}{Lemma}
\newtheorem*{conj*}{Conjecture}
\newtheorem*{ques*}{Question}
\newtheorem*{prob*}{Problem}
\newtheorem*{exam*}{Example} 
\newtheorem*{notn*}{Notation}
\theoremstyle{remark}
\newtheorem*{rmk*}{Remark}
\newtheorem*{recall*}{Recall}
\newtheorem*{case*}{Case}
\newtheorem*{desiderata*}{Desiderata}

%%%%% Multiple Choice
\newlist{inlinechoices}{enumerate*}{1}
\setlist[inlinechoices]{%
	label=\textbf{\alph*.},
	itemjoin={\hspace{6em}},
	itemsep=0pt
}

%%%%% Roman Enumerate
\setlist[enumerate]{label=\textbf{\roman*}.}

%%%%% Figures
\newcommand{\fig}[2][]{%
	\begin{figure}[!h]
		\centering
		\includegraphics[width=0.5\textwidth, height=0.35\textwidth, keepaspectratio, #1]{Figures/#2.png}
	\end{figure}
}

%%%%% Extra Credit
\NewDocumentCommand{\extracredit}{o+m}{%
	\stepcounter{question}%
	\item[\textbf{\thequestion.}]%
	\IfValueTF{#1}
		{(\textbf{Extra Credit}, #1 points)~}
		{(\textbf{Extra Credit})~}%
	#2
}

%%%%% Question Relabeling
\renewcommand{\questionlabel}{\textbf{\arabic{question}.}}
\renewcommand{\partlabel}{\textbf{\alph{partno}.}}

%%%%% Circle Letter
\newcommand\encircle[1]{%
	\tikz[baseline=(X.base)] 
	\node (X) [draw, shape=circle, inner sep=0] {\strut #1};
}

%%%%% Document Metadata
\usepackage{course-metadata} % Auto-loads coursemeta.tex from ancestor dirs

% --- Document-specific: set via \renewcommand in each quiz .tex file ---
\newcommand{\theQuizNumber}{\todo}
\newcommand{\theTitle}{Quiz \theQuizNumber}

% --- Course metadata: auto-populated from coursemeta.tex via course-metadata.sty ---
% These \the* aliases preserve backward compatibility with existing documents
% that use \renewcommand to override them; new documents need not touch them.
\newcommand{\theCourseNumber}{\GetCourseNumber}
\newcommand{\theCourseTitle}{\GetCourseTitle}
\newcommand{\theCourse}{\GetCourseSubject\ \GetCourseNumber\ \GetCourseTitle}
\newcommand{\theCourseSection}{\GetCourseSection}
\newcommand{\theSeason}{\GetSeason}
\newcommand{\theYear}{\GetYear}
\newcommand{\theTerm}{\GetTerm}
\newcommand{\theSchool}{\GetInstitution}

%%%%% Title Information
\title{\LARGE \theTitle}

% Heading and Footer
\pagestyle{headandfoot}
\firstpageheader{\theCourse}{}{Name: \underline{\hspace{2.5in}}}
\firstpageheadrule
\firstpagefooter{Section \theCourseSection}{}{\theSchool}
\firstpagefootrule
\runningheader{\theCourse}{}{Name: \underline{\hspace{2.5in}}}
\runningheadrule
\runningfootrule
\runningfooter{Section \theCourseSection}{}{\theSchool}
